\begin{abstract}
热风烘干过程中，药材表面首先受到环境作用，而内部温度和水分浓度以不同时间尺度向中心响应；当药材发生收缩时，传递距离和有效物性又会同步变化。为回答预热状态、全过程演化、全域达标时间及收缩影响四个递进问题，本文将药材理想化为轴对称长圆柱，构建统一的径向传热传质模型，并在结果层面区分局部状态、全域判停和移动表面三类输出。

针对问题一，采用附录2物性和固定半径的一维 Fourier 导热—Fick 扩散解耦模型，将附件1环境序列分段线性插值为连续边界；针对问题二，在相同几何下引入附录3的含水率相关密度、比热容、导热系数及温度—含水率相关扩散系数，形成同步更新的热质参数耦合模型；针对问题三，沿用问题二的全过程描述，以轴线、内部控制体和真实表面水分浓度的最大值定义连续终止事件；针对问题四，依据附件2半径轨迹和附录4物性，在材料坐标 $\xi=r/R(t)$ 下处理径向收缩，并对固定物理位置出域和实时表面分别输出。

空间上采用守恒有限体积法，内部界面物性取调和平均，表面状态由半网格关系重构；时间上采用 BDF 方法同步推进温度和水分。问题一计算表明，$1800\,\mathrm{s}$ 时中心与表面温度分别为 $33.5753$ 和 $36.7856\,{}^\circ\mathrm{C}$，对应水分浓度分别为 $2.5500$ 和 $1.5102\,\mathrm{kg/kg}$，热响应已进入内部而失水仍集中在表层。问题二的前三小时结果显示，温度场先于水分场趋于均匀，至 $3\,\mathrm{h}$ 轴线与表面温度分别为 $49.8495$ 和 $49.9664\,{}^\circ\mathrm{C}$，而水分浓度仍为 $1.7662$ 和 $1.0081\,\mathrm{kg/kg}$，内部传质滞后明显。

在固定半径、附录3物性条件下，问题三得到全域水分浓度首次达到 $0.15\,\mathrm{kg/kg}$ 的时间为 $57.17\,\mathrm{h}$。考虑附录4物性和附件2径向收缩后，问题四的连续临界时刻为 $182956.15\,\mathrm{s}$（约 $50.82\,\mathrm{h}$）；终点半径为 $1.2000\,\mathrm{cm}$，全域最大值位于轴线，表面水分浓度约为 $0.0525\,\mathrm{kg/kg}$。四工况对照表明，物性变化和几何收缩共同决定达标时间，不能将问题三与问题四的时间差全部归因于收缩。

网格加密、时间积分设置、全域守恒、跨问题对齐、零扩散极限、固定域退化复现及代表时刻边界根扫描共同支持给定假设下的数值可靠性与模型内部一致性。模型主要适用于圆柱主体的径向过程；长度不变、材料同比例收缩、环境数据末值延拓以及未显式加入蒸发潜热等处理，限定了结果向实际工艺推广时的适用范围。

\keywords{热风烘干\,径向传热传质\,变物性耦合\;移动边界\,全域达标判定}
\end{abstract}
